\scanpage{110}
\mspage{55}

dépendant du choix d'une orientation $e$ de $P_0$ -- de façon intrinsèque on a

\[
G_{0s}^+\simeq\Lambda\wedge_{\mu_2}\mu_0\tag*{(28)}
\]

\margin{le sens des deux symboles $\mu_2$ et $\mu_0$ est tout à fait différent !}

mais l'opération naturelle de $Z'$ sur $G_{0s}^+$, $\deleted{\unsure{\mathbb Z}}\Lambda_n$-module libre de rang un, se fait via un ``caractère''\ canonique

\[
Z'\xrightarrow{\chi_s}\Lambda_n^*,\tag*{(29)}
\]

\unsure{tu|On} peut songer à utiliser $\chi$ pour identifier $Z'$ à $\Lambda_n^*$ via à un ss-groupe de $\Lambda_n^*$. Mais \unsure{tantôt|bientôt|de fait|doit-il} on définira un isom can.

\[
Z'\simeq\Lambda_n^*/\pm1
\]

et $\chi_s$ sera défini par la formule

\[
\chi_s(\dot\lambda)=\lambda^2\qquad\text{pour }\lambda\in\Lambda_n^*\tag*{(30)}
\]

% slanted
\margin{cf formule $(*)$ p. $46'$}

qui montre qu'on ne doit pas s'attendre que (29) soit injectif -- donc il faut trouver un autre principe pour construire un iso. $Z'\simeq\Lambda_n^*/\pm1$.

\deleted{On suppose maintenant}

\[
\begin{aligned}
 &\deleted{(\unsure{\mathrm S|\delta}\mathrm T_0\,3)\quad\forall\lambda\in Z',\ s,t\in S,\ \text{on a}}\\
 &\deleted{\{s,t\}\in A_0\iff\{\lambda s,\lambda^{-1}t\}\in A_0\implies t\ne\lambda s}
\end{aligned}\tag*{\deleted{(31)}}
\]
% archive sheet number: 109
